% ===== supplementary/multi_agent_formalism.tex =====
% Archived from sec/04_trace_language.tex (multi-agent system definition).
% Removed from main paper per rewrite plan item 3.

\section{Multi-Agent Systems}
\label{sec:multi-archived}

\begin{definition}[Multi-agent system]
\label{def:multi}
A \emph{multi-agent system} is a finite tuple $\mathcal{M} =
(\mathcal{A}_1, \dots, \mathcal{A}_k, R)$ where each
$\mathcal{A}_i$ is an agent and $R \subseteq \Sigma^*\times\{1,
\dots,k\}\times\{1,\dots,k\}$ is an \emph{orchestration relation}
specifying which agent's emission, on which prefix, may be
consumed by which other agent.
\end{definition}

The trace language of $\mathcal{M}$ is the set of strings over the
common alphabet $\Sigma = \bigcup_i \Sigma_{\mathcal{A}_i}$ that
are obtainable by the orchestrated shuffle: pick an accepting run
of each $\mathcal{A}_i$, then interleave their traces in any order
consistent with $R$.  Formally,
\[
L(\mathcal{M}) \;=\;
\bigcup_{(\tau_1, \dots, \tau_k)\,\in\,
        L(\mathcal{A}_1)\times\cdots\times L(\mathcal{A}_k)}
\textsf{shuffle}_R(\tau_1, \dots, \tau_k)
\]
where $\textsf{shuffle}_R$ is the standard interleaving
restricted by $R$~\cite{hopcroft2006automata}.
